% TOP_PROBLEM: {"id": 30003556, "title": "Lehmer’s Mahler Measure Problem", "category_id": 1, "category": {"id": 1, "name": "number_theory", "display_name": "Number Theory", "description": "Properties of integers, prime numbers, Diophantine equations.", "slug": "number-theory", "order_index": 1, "created_at": "2026-07-31T15:26:25.670Z"}, "status": "open"}
% FIRST_POSTED: 2026-09-27
% Standalone research entry 206; batch 203--207.
% Original section material: Pasted text(3).txt, source edition 22 September 2026.
% Research additions: 27 September 2026.
% Compile twice with: lualatex 30003556.tex
% !TeX program = lualatex
% Compile twice: lualatex open_problems_500.tex
% All references and prose are embedded; no BibTeX database or external inputs.
\documentclass[11pt,a4paper]{article}
\usepackage[margin=24mm,headheight=14pt]{geometry}
\usepackage{fontspec}
\setmainfont{Latin Modern Roman}
\setsansfont{Latin Modern Sans}
\setmonofont{Latin Modern Mono}
\usepackage{amsmath,amssymb,mathtools}
\usepackage{unicode-math}
\setmathfont{Latin Modern Math}
\usepackage{microtype}
\usepackage{enumitem,xurl,needspace}
\usepackage[dvipsnames]{xcolor}
\usepackage{hyperref}
\hypersetup{unicode=true,colorlinks=true,linkcolor=MidnightBlue,urlcolor=MidnightBlue,
 pdftitle={500 Mathematical Problem Statements},pdfauthor={Source-annotated compilation},bookmarksnumbered=true}
\usepackage{bookmark}
\usepackage{tocloft}
\setlength{\cftsecnumwidth}{3em}
\cftsetpnumwidth{2.5em}
\cftsetrmarg{4em}
\usepackage{titlesec}
\titleformat{\section}{\Large\bfseries\raggedright}{\thesection}{0.7em}{}
\usepackage{fancyhdr}
\pagestyle{fancy}\fancyhf{}
\fancyhead[L]{\small\sffamily Mathematical problem statements}
\fancyhead[R]{\small Source edition 22}
\fancyfoot[C]{\thepage}
\setlength{\parindent}{0pt}
\setlength{\parskip}{5pt plus 1pt}
\setlength{\emergencystretch}{3em}
\setcounter{tocdepth}{1}
\setcounter{secnumdepth}{1}
\setlist[itemize]{leftmargin=3em,itemsep=5pt,topsep=4pt,parsep=0pt}
\allowdisplaybreaks
\newcommand{\N}{\mathbb{N}}
\newcommand{\Z}{\mathbb{Z}}
\newcommand{\Q}{\mathbb{Q}}
\newcommand{\R}{\mathbb{R}}
\newcommand{\C}{\mathbb{C}}
\newcommand{\F}{\mathbb{F}}
\newcommand{\A}{\mathbb{A}}
\newcommand{\PP}{\mathbb P}
\newcommand{\E}{\mathbb{E}}
\newcommand{\eps}{\varepsilon}
\newcommand{\dd}{\,\mathrm d}
\newcommand{\sref}[2]{\hyperlink{source:#1:#2}{[S#2]}}
\newcommand{\eref}[2]{\hyperlink{extra:#1:#2}{[E#2]}}
% Turn the next switch off for a shorter reading copy. The quotations preserve
% every exactTarget field, including qualifications omitted from a short summary.
\newif\ifshowcatalogue
\showcataloguetrue
\newcommand{\cataloguescope}[1]{\ifshowcatalogue
 \paragraph{Exact scope in the supplied catalogue.}
 \begin{quote}\small #1\end{quote}\fi}

% Standalone wrapper; no external inputs or bibliography files are required.
\fancyhead[L]{\small\sffamily Research notebook: section 30003556}
\fancyhead[R]{\small 27 September 2026}
\hypersetup{pdftitle={Research attempt 206},pdfauthor={Research notebook}}
\setlength{\parskip}{4pt plus 1pt}
\setlist[itemize]{before=\small\interlinepenalty10000,itemsep=3pt}
\begin{document}
\setcounter{section}{0}
% ================================================================
% problemId: 30003556
% Canonical problem ID: 30003556
% exactTarget SHA-256: 1e841796adcc455ab436f0853c9a9152f9178ea07b0a0ae85e4e487f75c43b31
\clearpage
\section*{\texorpdfstring{Lehmer’s Mahler Measure Problem}{Lehmer’s Mahler Measure Problem}}\label{30003556}
\subsection{Definitions and mathematical statement}
If a monic polynomial $P\in\Z[X]$ factors as
$P(X)=\prod_{j=1}^d(X-\alpha_j)$ over $\C$, its multiplicative Mahler measure is
\[
 M(P)=\prod_{j=1}^d\max\{1,|\alpha_j|\}.
\]
The selected gap assertion is that an absolute $\eps_0>0$ exists with
$M(P)>1\Longrightarrow M(P)\ge1+\eps_0$ for every monic integer polynomial of arbitrary degree. Equivalently, either establish a gap immediately above one or produce a sequence with $M(P_j)>1$ and $M(P_j)\to1$. No fixed-degree bound can substitute for an absolute bound across all degrees, and the polynomials need not be irreducible.
\par\smallskip\textit{Formulation and concept references:} \sref{30003556}{1}.
\cataloguescope{Prove that there exists an absolute epsilon > 0 such that every monic integer polynomial P with Mahler measure M(P) > 1 satisfies M(P) >= 1 + epsilon, or construct a sequence of such polynomials whose Mahler measures decrease to 1.}
\subsection{Short English statement}
Is there a universal gap above one among Mahler measures of monic integer polynomials, or can such measures approach one from above?
% BEGIN RESEARCH ADDITION 206
\subsection{Research attempt}

\paragraph{Current frontier.}
For unrestricted degree the established general lower bounds have
Dobrowolski-type dependence on the degree, while Smyth's theorem gives
a uniform gap for nonreciprocal algebraic integers of measure greater than one.\eref{30003556}{1}
The Schinzel--Zassenhaus theorem gives a lower bound for the \emph{house}
$\max_j|\alpha_j|$, not the required uniform lower bound for the product
$M(P)$. Dimitrov's Theorem 1 gives $\max_j|\alpha_j|\ge2^{1/(4d)}$
for monic irreducible noncyclotomic polynomials other than $X$;
the Annals publication page records acceptance in August 2026.\eref{30003556}{2}
For each fixed bound on the number of nonzero monomials, there is a gap
above measure one, independent of the exponents; that bound cannot be
held fixed in a hypothetical counterexample sequence.\eref{30003556}{3}
None of these statements substitutes for the selected degree-independent
assertion over all monic integer polynomials.

\paragraph{Attack route.}
Resultant methods would need arithmetic amplification that does not
deteriorate with degree. A sparse-polynomial construction cannot work
with bounded support size, by the preceding theorem.\eref{30003556}{3}
The route developed here combines exact integer power sums with positive
semidefinite moment matrices of points on the unit circle. Reciprocal
pairing causes a \emph{quadratic}, rather than linear, radial error.
This produces explicit lower-bound certificates and a normalization
that survives the substitution $X\mapsto X^N$. The unresolved question
is whether certificates of uniformly positive normalized strength always
exist, not merely whether some moment eventually detects an off-circle
root.

\paragraph{Attempt.}
\textbf{1. Reduction and normalization.}
Write $m(P)=\log M(P)$. Multiplicativity of $M$ follows immediately by
concatenating roots. A monic integer polynomial with $M(P)>1$ therefore
has a monic irreducible factor $Q$ with $1<M(Q)\le M(P)$. Factors $X$
can be removed. If $|Q(0)|\ge2$, then
$M(Q)\ge|Q(0)|\ge2$. The unresolved small-measure regime is thus
algebraic units. Smyth's nonreciprocal bound further reduces a prospective
sequence tending to one to reciprocal units.\eref{30003556}{1}
For an irreducible reciprocal unit of measure greater than one we can
use $Q(X)=X^dQ(1/X)$ and $Q(0)=1$: the anti-reciprocal possibility
would have a root at $1$ and hence could not be such an irreducible
polynomial. The estimates below also apply to reducible reciprocal
monic polynomials with constant term one.

Fix such a polynomial $P$, with roots counted with multiplicity, and set
\[
 \alpha_j=r_je^{i\theta_j},\quad
 m=\sum_{r_j>1}\log r_j,\quad
 h=\max_{r_j>1}\log r_j>0,
 \qquad
 s_k=\sum_{j=1}^d\alpha_j^k\in\Z\quad(k\ge0).
\]
Put $s_0=d$ and extend by $s_{-k}=s_k$. Reciprocal symmetry justifies
this extension as the actual negative power sum; reality of the
coefficients makes these sums real. Define the projected phase moments
\[
 u_k=\sum_{j=1}^de^{ik\theta_j},\qquad u_{-k}=u_k.
\]
The phases occur in conjugate pairs, so $u_k$ is real.

\textbf{Lemma 206.1 (quadratic reciprocal radial error).}
For every integer $k$,
\[
 |s_k-u_k|\le k^2 h m\cosh(|k|h). \tag{206.1}
\]

\emph{Proof.}
An outside root $\alpha=e^{\ell+i\theta}$ is paired with the inside
root $1/\overline\alpha=e^{-\ell+i\theta}$; the polynomial is both
real and reciprocal, so this pairing respects multiplicities. Their
contribution to $s_k-u_k$ is
\[
 \bigl(e^{k\ell}+e^{-k\ell}-2\bigr)e^{ik\theta}
 =2(\cosh(k\ell)-1)e^{ik\theta}.
\]
Unit-circle roots contribute zero. Since
$2(\cosh t-1)\le t^2\cosh t$ for real $t$, summing over the outside
roots gives
\[
 |s_k-u_k|
 \le k^2\cosh(|k|h)\sum_{r_j>1}(\log r_j)^2
 \le k^2hm\cosh(|k|h).
\]
The elementary hyperbolic-cosine inequality follows by integrating
$\cosh$ twice from zero, or termwise from its power series. All sums
over roots are finite.\hfill$\square$

The quadratic cancellation would be lost by estimating each root and
its reciprocal separately. Importantly $h$, the largest individual
radial deviation, can be much smaller than the total deviation $m$.
Replacing it by $m$ prematurely destroys useful substitution invariance.

\textbf{2. Exact integer certificates from moment positivity.}
Let $a_0,\ldots,a_L$ be distinct nonnegative integer indices and
$v=(v_0,\ldots,v_L)\in\Z^{L+1}$. Define
\[
 T=(s_{a_i-a_j})_{0\le i,j\le L},\qquad
 G=(u_{a_i-a_j})_{0\le i,j\le L}.
\]
The second matrix is positive semidefinite because
\[
 v^T Gv=\sum_{j=1}^d
 \left|\sum_{i=0}^L v_i e^{ia_i\theta_j}\right|^2\ge0.
\]
If the exact integer calculation gives $v^TTv=-q<0$, Lemma 206.1
therefore yields
\[
 q\le hm\,H_h(a,v),\qquad
 H_h(a,v)=\sum_{i,j=0}^L|v_iv_j|(a_i-a_j)^2
                   \cosh(|a_i-a_j|h).
\]
In particular,
\[
 \boxed{\displaystyle m(P)\ge\frac{q}{h\,H_h(a,v)}.}
 \tag{206.2}
\]
The denominator is positive for a negative certificate, since otherwise
$v$ has at most one nonzero entry and $v^TTv=d\|v\|_2^2\ge0$.
If a rigorous upper bound $\bar h\ge h$ is available, the weaker but
certified estimate $m(P)\ge q/(\bar h H_{\bar h}(a,v))$ follows by
monotonicity. Approximate eigenvectors may suggest $v$, but the sign of
$v^TTv$ must be checked over the integers.

There is always \emph{some} negative certificate when $m(P)>0$.
Otherwise every $2\times2$ matrix with indices $0,k$ is positive
semidefinite, giving $|s_k|\le d$ for all $k$. The power series
\[
 \sum_{k\ge0}s_kz^k=\sum_{j=1}^d\frac{1}{1-\alpha_jz}
 \tag{206.3}
\]
would then be analytic for $|z|<1$. An outside root produces a pole
at $z=1/\alpha_j$ in that disk. Repeated copies of this root add
identical nonzero residues; distinct roots have different pole locations,
so the pole cannot cancel. This contradicts (206.3), initially valid
near zero and hence throughout its domain by analytic continuation.
Thus some $|s_k|>d$, and $v=(1,-\operatorname{sgn}s_k)$ at indices
$0,k$ is negative. This proves existence but gives no uniform strength
for (206.2).

\textbf{3. A substitution stress test, and why the normalization helps.}
For any monic polynomial $Q$ and positive integer $N$,
\[
 M(Q(X^N))=M(Q),\qquad
 \operatorname{house}(Q(X^N))=
 \operatorname{house}(Q)^{1/N}. \tag{206.4}
\]
Indeed the $N$ roots above a nonzero root $\beta$ all have modulus
$|\beta|^{1/N}$, and their factors in the product defining $M$ multiply
to $\max(1,|\beta|)$. Zeros also obey the formula. Consequently small
house alone cannot manufacture a Lehmer counterexample.

Moreover the power sums of $Q(X^N)$ are zero when $N$ does not divide
$k$, and equal $N s_{k/N}(Q)$ when it does. Thus any fixed finite prefix
of positive-index traces can be identically zero for a polynomial of
measure greater than one. The cyclotomic polynomial product
$X^{N\deg Q}+1$ has the same degree and the same prefix. A test using
only a fixed initial trace window therefore cannot decide the original
problem. Substitution may make a polynomial reducible, which is allowed
by the original target; this example does not by itself refute every
possible irreducibility-sensitive test.

Nevertheless (206.2) passes a more demanding check. For reciprocal $Q$,
replace $a_i$ by $Na_i$ and keep $v$ when passing to $Q(X^N)$. Then
\[
 q_N=Nq,\qquad h_N=h/N,\qquad
 H_{h_N}(Na,v)=N^2H_h(a,v).
\]
Hence the certificate value $q_N/(h_NH_{h_N}(Na,v))$ is exactly the
original value. This is why keeping the quadratic error and the separate
quantity $h$ is preferable to a crude bound involving only $km$.
The method does not automatically lose its estimate merely because
all exponents have been multiplied by $N$.

\textbf{4. A test certificate with exact verification.}
Take the classical degree-ten test polynomial
\[
 P_\mathrm L(X)=X^{10}+X^9-X^7-X^6-X^5-X^4-X^3+X+1.
\]
Its role as Lehmer's example is historical; the calculations that follow
are performed directly.\eref{30003556}{1}
Newton's identities give, without root approximations,
\[
 (s_0,\ldots,s_8)=(10,-1,1,2,1,4,4,6,1).
\]
With consecutive indices $a_i=i$, choose
\[
 v=(2,3,2,1,0,-1,-2,-3,-2).
\]
Exact integer multiplication gives
\[
 v^TTv=-106,\qquad
 \sum_{i,j=0}^8|v_iv_j|(i-j)^2=4352. \tag{206.5}
\]
The vector was found by rounding a numerically computed negative
eigenvector. Only the integer identities (206.5), not the floating-point
eigenvalue, enter the certificate.

For a rigorous radial bound, use
\[
 P_\mathrm L(X)=X^5 Q(X+X^{-1}),\qquad
 Q(t)=t^5+t^4-5t^3-5t^2+4t+3.
\]
A Sturm sequence, with each member independently scaled by a positive
constant, is
\[
\begin{gathered}
 Q(t),\quad 5t^4+4t^3-15t^2-10t+4,\\
 54t^3+60t^2-90t-71,\quad96t^2+117t-38,
 \quad1167t+1390,\quad1.
\end{gathered}
\]
Put $a=117628/100000$, $b=117629/100000$ and
$t_a=a+a^{-1}$, $t_b=b+b^{-1}$. Direct rational sign evaluation gives
\[
\begin{array}{c|rrrrrr}
 t&-\infty&-2&2&t_a&t_b&+\infty\\\hline
 \text{sign variations}&5&5&1&1&0&0.
\end{array}
\]
Thus four roots of $Q$ lie in $(-2,2)$ and its remaining root lies in
$(t_a,t_b)\subset(2,\infty)$. The corresponding roots of $P_\mathrm L$
are eight unit-circle roots and a positive pair $\lambda,\lambda^{-1}$
with $a<\lambda<b$. Since
\[
 e^{1/6}>1+\frac16+\frac1{72}+\frac1{1296}>b,
\]
we obtain $h=m(P_\mathrm L)=\log\lambda<1/6$.
Finally $\cosh(4/3)<21/10$: sum its terms through degree ten and bound
the tail beginning with $(4/3)^{12}/12!$ by the geometric ratio
$(4/3)^2/(14\cdot13)$; the resulting rational upper bound is
$6491958743731/3200164983675<21/10$.
Equations (206.2)--(206.5) now prove the explicit, albeit weak, bound
\[
 m(P_\mathrm L)>\frac{106}{(1/6)(21/10)4352}
 =\frac{265}{3808}. \tag{206.6}
\]
This is a check of the certificate mechanism on a known example, not
an improved known lower bound for that polynomial or for the problem.
Numerical evaluation gives $\lambda\approx1.1762808182599175$ and
$m\approx0.1623576120077381$; using the actual $h$ in (206.2) gives a
certificate value approximately $0.09753$. Those decimals are numerical
illustrations only; (206.6) uses rational and exact integer verification.

\textbf{5. An infinite construction that fails for an exact reason.}
The candidate family $X^{2n}+aX^n+1$, $a\in\Z$, has
\[
 M(X^{2n}+aX^n+1)=
 \begin{cases}
 1,&|a|\le2,\\[2pt]
 (|a|+\sqrt{a^2-4})/2,&|a|\ge3.
 \end{cases}
\]
Solve $Y^2+aY+1=0$ and use (206.4). Thus taking $n\to\infty$ can
make every root close to the unit circle without decreasing the Mahler
measure at all. The smallest measure greater than one in this entire
family is $(3+\sqrt5)/2$. More generally, any hypothetical sequence
with $M\downarrow1$ must have unbounded monomial count, by the established
bounded-support theorem.\eref{30003556}{3}

\paragraph{Stress test.}
Reciprocity, including its interaction with complex conjugation, is
essential to (206.1); for a general unit the linear radial error need
not cancel. The phase moment matrix is positive semidefinite for an
explicit reason, whereas the integer trace matrix need not be. The
Sturm calculation avoids assuming from a plot that all the other roots
lie on the unit circle. Existence of some negative certificate follows
from a pole argument, not from extrapolating the finite eigenvalue search.

The unresolved uniform step is precise. It would suffice to find an
absolute $c>0$ such that every remaining irreducible reciprocal $P$
with $M(P)>1$ admits integer indices and a vector with
\[
 -v^TTv\ \ge\ c\,h H_h(a,v)>0. \tag{206.7}
\]
Then (206.2), together with the established reductions, would give the
requested absolute gap. No bound of the form (206.7) has been proved.
An integer negative quadratic form is at most $-1$, but its denominator
in (206.2) can grow; integrality alone cannot be used to erase that
growth. A proof of (206.7) would settle Lehmer's problem and so cannot
be regarded as an innocuous numerical conditioning assertion.

\Needspace{10\baselineskip}
\paragraph{Outcome.}
\textbf{Outcome: Exploratory approach with unresolved gap.}

The attempt derives a quadratic reciprocal-root estimate and an exact
integer moment-certificate inequality, proves that some negative
certificate always exists off the unit circle, and checks the method
with an independently verifiable degree-ten certificate. Its normalized
bound is invariant under power substitution, removing one concrete
failure mode of house-based arguments. No claim is made that these
notebook lemmas improve the established general frontier or are new in
the literature.

A complete proof requires uniformly strong certificates such as (206.7),
or another argument preventing their strength from tending to zero
along increasing-degree, increasing-support polynomials. The explicit
families tested here do not yield a counterexample sequence, and the
finite computation does not exclude such sequences in general.

% Keep the sources heading with a useful initial bibliography block.
\Needspace{8\baselineskip}
% END RESEARCH ADDITION 206
\subsection{Sources}
\begin{itemize}\raggedright
\item[\textnormal{[S1]}]\hypertarget{source:30003556:1}{} Lehmer's Mahler Measure Problem (MathWorld).
\url{https://mathworld.wolfram.com/LehmersMahlerMeasureProblem.html}.
{\small\textit{Catalogue formulation source.}}
% BEGIN REFERENCE ADDITION 206
\item[\textnormal{[E1]}]\hypertarget{extra:30003556:1}{}
C. Smyth, \emph{The Mahler measure of algebraic numbers: a survey},
in \emph{Number Theory and Polynomials}, London Mathematical Society
Lecture Note Series 352, Cambridge University Press (2008), 322--349;
arXiv:math/0701397v2. Sections 1, 4.2 and 5 discuss Lehmer's polynomial,
Dobrowolski's bound and the nonreciprocal theorem.
\url{https://arxiv.org/abs/math/0701397}.
\item[\textnormal{[E2]}]\hypertarget{extra:30003556:2}{}
V. Dimitrov, \emph{A proof of the Schinzel--Zassenhaus conjecture on
polynomials}, arXiv:1912.12545, Theorem 1; accepted version titled
\emph{The Schinzel--Zassenhaus conjecture}, Annals of Mathematics,
forthcoming. The journal page records acceptance on 27 August 2026;
consulted 27 September 2026. This is a house theorem, not a solution
of the Mahler-measure gap problem.
\url{https://arxiv.org/abs/1912.12545}.
\url{https://annals.math.princeton.edu/articles/22975}.
\item[\textnormal{[E3]}]\hypertarget{extra:30003556:3}{}
E. Dobrowolski and C. Smyth, \emph{Mahler measures of polynomials that
are sums of a bounded number of monomials}, arXiv:1606.04376 (2016),
Theorem 2 and its proof in Section 4: isolation of zero in the set
of logarithmic measures for bounded monomial count.
\url{https://arxiv.org/abs/1606.04376}.
% END REFERENCE ADDITION 206
\end{itemize}

\end{document}
